\section{Geometric identities and profile consequences}\label{sec:geometry}

We now extract three consequences needed in the curvature estimates:
uniform bounds for the warping slopes, a strict inner separation from the
cone, and small Ricci curvature on fixed parabolic annuli. We distinguish
time differentiation at a fixed point of the manifold from differentiation
at a fixed value of the evolving radius \(r\).

\subsection{Warping equations}

Let \(x\) be a radial coordinate fixed on the manifold, so that
\(ds=\chi(x,t)\,dx\). Subscripts \(s\) denote arclength differentiation
at the indicated time. On a regular orbit the sectional curvatures are
\begin{equation}\label{eq:sectional}
\begin{gathered}
h=-\phi_{ss}/\phi,\qquad h_q=-r_{ss}/r,\qquad
j=(1-\phi_s^2)/\phi^2,\\
\ell=(1-r_s^2)/r^2,\qquad
\mu=-\phi_s r_s/(\phi r).
\end{gathered}
\end{equation}
They correspond, respectively, to radial--\(\Sp^p\), radial--\(\Sp^q\),
\(\Sp^p\)-tangent, \(\Sp^q\)-tangent, and mixed planes. These formulas
also follow directly by computing the connection of the warped metric.
The Ricci eigenvalues in the radial and the two fiber directions are
\begin{equation}\label{eq:ricci-eigenvalues}
\lambda_0=ph+qh_q,\qquad
\lambda_p=h+(p-1)j+q\mu,\qquad
\lambda_q=h_q+(q-1)\ell+p\mu.
\end{equation}
Consequently
\begin{equation}\label{eq:warping-flow}
\begin{aligned}
\dt\big|_x r&=r_{ss}+(q-1)(r_s^2-1)/r+p\phi_s r_s/\phi,\\
\dt\big|_x\phi&=\phi_{ss}+(p-1)(\phi_s^2-1)/\phi+q r_s\phi_s/r,\\
[\dt\big|_x,\partial_s]&=\lambda_0\partial_s.
\end{aligned}
\end{equation}
These are also the usual doubly warped Ricci-flow equations
in~\cite[Section~2]{Stolarski2019}.

\begin{lemma}[Global elementary bounds]\label{lem:gradients}
For each flow in Proposition~\ref{prop:input}, there is a time-independent
constant \(C\) such that
\[
r\le C,\qquad |r_s|+|\phi_s|\le C.
\]
In particular, \(r\) is uniformly Lipschitz in the metric \(g(t)\).
Also \(R\ge-C\) throughout the flow.
\end{lemma}

\begin{proof}
At a positive maximum of \(r\), its first derivative vanishes and
\eqref{eq:warping-flow} gives \(r_t\le-(q-1)/r\).
Differentiating the first equation with the stated commutator gives
\begin{equation}\label{eq:slope-flow}
\begin{split}
\dt\big|_x r_s={}&(r_s)_{ss}
 +\left(p\frac{\phi_s}{\phi}+(q-2)\frac{r_s}{r}\right)(r_s)_s\\
&+\left[\frac{(q-1)(1-r_s^2)}{r^2}
       -p\frac{\phi_s^2}{\phi^2}\right]r_s.
\end{split}
\end{equation}
At a positive interior maximum with \(r_s>1\), the reaction has negative
sign; the corresponding sign is positive at a negative minimum below
\(-1\). Thus \(|r_s|\) is bounded by the larger of one and its initial
supremum. Interchanging \((r,q)\) and \((\phi,p)\) proves the same bound
for \(|\phi_s|\). On the complete interval between poles, the endpoint
values are \(r_s=1,-1\) and \(\phi_s=0\), so every offending extremum is
interior. The argument is applied first on a compact time slab; the
singular radial coefficients at the endpoints cause no extra boundary
condition. Finally
\[
(\dt-\Delta)R=2|\Ric|^2\ge0
\]
and the compact maximum principle imply
\(R(\cdot,t)\ge\min_M R(\cdot,0)\).
\end{proof}

Where \(r_s>0\), regard \(u=\log\phi\) and \(v=r_s=\sqrt z\) as functions
of \((r,t)\). From now on their time derivatives fix \(r\).
Subtracting the advection \((\dt|_x r)\partial_r\) in
\eqref{eq:warping-flow} gives
\begin{equation}\label{eq:sideways}
\begin{aligned}
u_t&=zu_{rr}+\frac{q-1+z}{r}u_r-(p-1)e^{-2u},\\
v_t&=zv_{rr}+\frac{q-1-z}{r}v_r
       +\frac{(q-1)v(1-v^2)}{r^2}-pv^3u_r^2.
\end{aligned}
\end{equation}
For later use, the geometric heat operator on a scalar composition with
\(r>0\) is
\begin{equation}\label{eq:operator}
\D F:= (\dt-\Delta)F(r,t)
=F_t|_r-zF_{rr}-\frac{q-1+z}{r}F_r.
\end{equation}
Indeed \(\Delta r=r_{ss}+q r_s^2/r+p\phi_s r_s/\phi\), so
\(r_t-\Delta r=-(q-1+z)/r\).
Equations~\eqref{eq:sideways} are written in a monotone radius coordinate.
The composition identity~\eqref{eq:operator} is geometric and remains valid
across the equator. The same warping equations give
the global identity, off the poles,
\begin{equation}\label{eq:ratio-heat}
(\dt-\Delta)\log(\phi/r)=\frac{q-1}{r^2}-\frac{p-1}{\phi^2}.
\end{equation}

\subsection{Consequences of the construction}

The first estimate is Stolarski's strict cone separation
\cite[Proposition~5.1(2)]{Stolarski2019}, expressed in the present scales.
We include its short proof to record the fixed positive gap as the cap
shrinks. Its constant may be small and may depend on the chosen flow;
it will be used only after the dimension and mode have been fixed.

\begin{lemma}[A strict inner gap]\label{lem:gap}
After choosing the tolerances sufficiently small, there are constants
\(\varepsilon_1>0\) and \(C<\infty\) such that, at all sufficiently late
times,
\begin{equation}\label{eq:strict-gap}
\phi/r\ge A/B+\varepsilon_1\quad(0<r\le Y_1\theta),
\qquad \phi(Y_1\theta,t)\le C\theta.
\end{equation}
\end{lemma}

\begin{proof}
At \(\gamma_1=Y_1\delta^{\nu/d}\),
\(\delta^\nu\gamma_1^{-d}=Y_1^{-d}\), while the relative error in the
leading small-\(\gamma\) asymptotic of \(U_k\) tends to zero.
Choose \(\eta_0^U<c_1/4\). Equation
\eqref{eq:profile-bounds} then gives
\[
0<\tfrac12c_1Y_1^{-d}\le\Ut(\gamma_1,\tau)\le C_{q,k,Y_1}
\]
at late times. This proves both assertions at the interface. Since
\(\partial_{\log r}\log(\phi/r)=f-1\le0\), the lower bound propagates
inward. For example, a fixed positive number below
\((A/B)(\exp(c_1Y_1^{-d}/2)-1)\) is admissible as \(\varepsilon_1\).
\end{proof}

\begin{lemma}[Bounds used in the dimension estimate]\label{lem:uniform}
The construction parameters can be chosen so that, for all sufficiently
large \(q\) and all sufficiently late times,
\begin{equation}\label{eq:uniform}
|z-1|\le C_0/q,\qquad r|u_r|\le2,\qquad
(p-1)r^2e^{-2u}\le(q-1)(1+C_0/q)
\end{equation}
on \(0<r\le\sqrt\delta\), where \(C_0\) is a numerical constant
independent of \(q\) and \(k\). This interval is below the equator.
\end{lemma}

\begin{proof}
For \(0<\gamma\le1\), the fixed profiles and their errors satisfy
\[
|\Ut|+\gamma|\Ut_\gamma|
\le C_{q,k}\delta^\nu\gamma^{-d},\qquad
|\Zt|\le C'_{q,k}\delta^\nu\gamma^{-d}
\]
above their respective lower cutoffs. Enlarge \(Y_1\), then \(Y_2\), until
the right sides at those cutoffs are at most \(1/q\).
They only decrease as \(\gamma\) increases. In the complementary inner
ranges use~\eqref{eq:inner-inputs}. Since
\[
1-B^2=\frac{2}{q+1},\qquad
ru_r=1+\gamma\Ut_\gamma,\qquad
(p-1)r^2e^{-2u}=(q-1)e^{-2\Ut},
\]
we obtain \(|z-1|\le3/q\), \(r|u_r|\le1+1/q\), and the last bound with
\(e^{2/q}\le1+4/q\) for large \(q\). Thus \(C_0=4\) suffices.
All dimension- or mode-dependent profile constants have been absorbed by
the cutoff choices; they do not enter \(C_0\).

The lower cutoffs tend to zero and the upper cutoff exceeds one at late
times. If the equator occurred before \(\gamma=1\), the same estimate up
to that endpoint would contradict its value \(z=0\), by continuity.
Thus the stated interval is available in the \(r\) coordinate.
\end{proof}

\begin{lemma}[Global ratio and parabolic annuli]\label{lem:annuli}
For a flow chosen as above, \(\phi/r\ge c_0>0\) globally off the poles.
For every fixed \(0<a_1<a_2<\infty\), at all sufficiently late times,
\begin{equation}\label{eq:annular-ricci}
|\Ric|\le C_{a_1,a_2}\delta^{\nu-1}
\qquad(a_1\sqrt\delta\le r\le a_2\sqrt\delta).
\end{equation}
The constants and the required starting time may depend on the annulus.
\end{lemma}

\begin{proof}
The positive large-\(\gamma\) coefficient of \(U_k\), with a tolerance
small compared with that coefficient, gives \(\Ut\ge0\) for all
\(\gamma\ge\Gamma_0\) in the profile region, where \(\Gamma_0\) is a
fixed sufficiently large number. At the initial time use the initial
profile bounds up to \(G_{\mathrm{init}}\), followed by the outer sign
condition beyond that same endpoint. Increase the construction's
starting rescaled time so that \(\Gamma_0\) belongs to the controlled
region for every subsequent time. On the full exterior
\(\{r\ge\Gamma_0\sqrt\delta\}\), compare \(\log(\phi/r)\) with
\(\log(A/B)\) in~\eqref{eq:ratio-heat}. The latter is an exact equilibrium
value of its reaction. Initial and moving lateral data have the correct
sign. On each compact time slab the reaction is locally Lipschitz there,
so comparison applies. Both hemispheres are included, and the equator is
not a boundary. Thus \(\phi/r\ge A/B\) on the exterior.
On \(1\le\gamma\le\Gamma_0\), profile control bounds \(\Ut\) below;
on \(\gamma\le1\) use Lemma~\ref{lem:uniform}.
Earlier compact time intervals have a positive lower ratio by smoothness,
positivity of \(\phi\), and boundedness of \(r\).

For a fixed positive compact \(\gamma\)-interval, the pulled-back metric
\(\delta^{-1}g\) is
\[
\frac{d\gamma^2}{z}+(A/B)^2\gamma^2e^{2\Ut}g_{\Sp^p}
+\gamma^2g_{\Sp^q}.
\]
It differs in \(C^2\) by \(O(\delta^\nu)\) from the metric with
\(z=B^2\), \(\Ut=0\). The latter is the Ricci-flat cone with link
\(A^2g_{\Sp^p}+B^2g_{\Sp^q}\). Curvature depends smoothly on a positive
metric and its first two derivatives on this fixed annulus. Hence its
Ricci norm is \(O(\delta^\nu)\) in the rescaled metric. Scaling back gives
\eqref{eq:annular-ricci}.
Any prescribed fixed annulus eventually lies inside the profile region;
as above, positivity of \(z\) there excludes the equatorial endpoint.
\end{proof}
